% Setup fonts -- see fontspec/mathspec documentation.
% Fonts are loaded from .ttf and .otf files in the core/fonts/ directory. We NEVER use system fonts.

% Change allocator to allow more than 16 alphabets per document
\makeatletter
    \def\new@mathgroup{\alloc@8\mathgroup\mathchardef\@cclvi}
    \patchcmd{\document@select@group}{\sixt@@n}{\@cclvi}{}{}
    \patchcmd{\select@group}{\sixt@@n}{\@cclvi}{}{}
\makeatother

\defaultfontfeatures{
    Mapping         = tex-text,
    Ligatures       = TeX
}

\setmainfont[
    Path            = assets/fonts/MinionPro/ ,
    Extension       = .otf ,
    UprightFont     = *-Regular ,
    BoldFont        = *-Bold ,
    ItalicFont      = *-It ,
    BoldItalicFont  = *-BoldIt
]{MinionPro}

% The monospace font takes no `tex-text` mapping, and this is not cosmetic. That mapping turns
% a straight `"` into `”`, a `'` into `’` and `--` into an en dash, which is exactly right for prose
% and wrong in every code listing: a solution's Python printed `print(f”Najvacsiu...”)`, which will
% not run, and neither the source nor the TeX showed it -- the substitution happens in the font.
% An empty `Mapping` overrides the default above for this family alone; the ligature key stays,
% because UbuntuMono has no f-ligatures to suppress.
\setmonofont[
    Path            = assets/fonts/UbuntuMono/ ,
    Extension       = .ttf ,
    UprightFont     = *-r ,
    BoldFont        = *-b ,
    ItalicFont      = *-i ,
    BoldItalicFont  = *-bi ,
    Mapping         =
]{UbuntuMono}

\setsansfont[
    Path            = assets/fonts/Verdana/ ,
    Extension       = .ttf ,
    UprightFont     = *-r ,
    BoldFont        = *-b ,
    ItalicFont      = *-i ,
    BoldItalicFont  = *-bi
]{verdana}

\DeclareFontFamily{U}{skulls}{}
\DeclareFontShape{U}{skulls}{m}{n}{ <-> skull }{}
\NewDocumentCommand{\skull}{}{\text{\usefont{U}{skulls}{m}{n}\symbol{'101}}}
